\section{Discussion and limitations}
\label{sec:discussion}

\subsection{Contribution in relation to existing work}

The contribution described here is a specific correspondence between a planning method and its software representation. The Research Map relates a problem, objectives, literature and development topics, and an objective--method matrix; the inspected application makes these units editable, records proposals and confirmations separately, and calls attention to later units when an earlier one changes. Its agent assists in formulating and revising content within that structure. The constructed walkthrough shows how the correspondence can be read, rather than how successfully people use it. This is a design and implementation account, not a claim that the Research Map originated staged planning, human review, provenance, or AI-supported research ideation.

The closest methodological papers establish why those broad features cannot serve as novelty claims. Research LLM describes a human-in-the-loop research companion across broader phases \cite{guerard2026researchllm}; Wang et al. organize AI assistance for clinical question formulation in five stages \cite{wang2026fivestage}. The present article instead specifies the Research Map's own units and their interdependencies, including the companion matrix, and examines their translation into application states and checks. CoQuest concentrates on co-creating research questions and reports a researcher study \cite{liu2024coquest}; ResearchAgent and the biomedical system of Aamer et al. explore agentic idea development, critique, knowledge structures, or memory \cite{baek2025researchagent,aamer2026biomedical}. Their empirical findings cannot be transferred to this implementation. Conversely, the present system has no dedicated research-questions unit, so its coverage of that task is narrower than a label such as ``research-question agent'' would suggest.

Workflow-history and provenance work provides another necessary boundary. VisTrails and RaivenTracks make revisions or branches of computational work inspectable, while PROV-DM supplies a general model for describing provenance \cite{freire2006vistrails,hugie2026raiventracks,moreau2013prov}. The Research Map application retains revisions and explicit planning links, but its stale marking follows predefined element types. It neither computes exact impact along each stored link nor demonstrates compliance with PROV-DM. Public descriptions of planning products, particularly ProjectReady AI and Researchly AI, also report overlapping planning or revision assistance \cite{projectready2026guide,researchly2026site}. As those descriptions were not tested in authenticated workflows for this article, the comparison supports a narrow positioning of the documented method--system mapping, not a claim of product priority or feature exclusivity.

The application also exposes a useful distinction within ``human supervision.'' The researcher can inspect and confirm proposed academic content; in student authoring mode, a submitted stage enters a pending review that an authenticated advisor can approve or return with comments. The agent's suggestion, the author's confirmation, the advisor's workflow decision, and scientific validation are different acts. An attempted e-mail notification helps bring a pending review to the advisor's attention, but delivery or reading is not established by the inspected route. This distinction matters when interpreting the flowchart in Section~\ref{sec:system}: its arrows specify possible transitions and responsibilities, not observed use or proof of research quality.

\subsection{Limits of the implementation and evidence}

Several method-to-system correspondences remain partial. The detailed method specification includes research questions among items to validate, whereas the inspected schema places question wording in problem candidates without a separate research-questions unit \cite{franca2026detail}. Topic coverage and trace links record declared relationships; lexical alignment and field-completeness checks cannot decide whether a question is worthwhile, a source supports a claim, or a method can answer an objective. The present suggestion library is static rather than a demonstrated growing methodological knowledge base. Bibliographic identifiers and provider metadata aid identification, but do not verify the relevance or evidentiary support of a cited work.

Revision handling is similarly bounded. Confirming changed content can mark downstream element \emph{types} stale and preserve prior versions, thereby prompting reconsideration. This conservative rule may call for review of unaffected items and cannot prove that every substantively affected relation has been found. At finalization, the route uses an advisory completion predicate and omits blocking findings from its persisted summary. A completed or exported map must therefore not be read as a certificate of coherence. Earlier stage checks and student--advisor approval can constrain transitions, but neither substitutes for disciplinary judgment about the research plan.

The evidence in this article consists of the source method documents, a directed comparison with published work and public product descriptions, inspection of the relevant implementation and local technical checks, and one constructed illustration. The checks support selected statements about program behavior at the examined layers; they are not a complete end-to-end test of every database, notification, agent, and export path. The fictional library case was not entered into the application and contains no generated agent output or observed advisor review. There are no participants, measured completion times, methodological ratings by the authors, or comparisons with an alternative planning process. Consequently, this article cannot establish that the system improves research quality, reduces delays, saves time, or is usable in routine supervision. Its claims concern documented design and inspected capabilities at the examined revision.

\subsection{Implications for further work}

The immediate design implication is to keep proposal, confirmation, review, and evidence status visible as separate information. More precise dependency tracking, explicit handling of unresolved final findings, verification of bibliographic support, and observable notification delivery would make the implementation's claims easier to inspect. These are possible extensions, not features demonstrated here. A separate empirical study could examine how researchers and advisors understand the map and its warnings, and whether the workflow changes planning decisions or effort. Such a study would require its own protocol and data; it is outside the evidentiary scope of this method-and-system article.
